\scanpage{224}
% continuation and framed definition
que \deleted{\unsure{pour}} sur
\[
 \boxed{\Phi=f_*(Z\underline{\Omega}^{1}_{X/S})/d f_*(\underline{\mathcal O}_X)}
 \tag*{(6)}
\]
, il y a une connexion intégrable naturelle, grâce à la formule de Cartan --- Les $\omega_i$ définissent des sections $\omega'_i$ de $\Phi$, et l'on peut former les $D_i\deleted{\unsure{\omega}}\omega'_i$. Je dis qu'il suffit qu'on ait
% slanted
\margin{attention, ne marche tel quel que si $X\longrightarrow S$ est steinien\unsure{!|?}}
\[
 \sum D_i\omega'_i=0.\tag*{(7)}
\]

% calculation using a connection
Lorsque $X/S$ est muni d'une connexion intégrable, alors le calcul de $\sum D_i\omega'_i$ se fait ainsi : \inserted{$f_*(\Omega^*_{X/S})$ sera aussi stratifié, donc} chaque $D_i$ se remonte \deleted{par la stratification en} \inserted{en un opérateur} $\overline D_i$ \deleted{sur} $\underline{\Omega}^*_X$ \inserted{commutant à la diff. ext. $d_{X/S}$} et \deleted{\unsure{comme} $\omega_i$ est une forme fermée, on peut former}

% entire enclosed block cancelled by several large oblique crosses
\deleted{\[
\boxed{\begin{aligned}
 (8)\quad &\alpha_{\deleted{\unsure{D_i}}\overline D_i}\omega_i
 =\alpha_{\overline D_i^0}\omega_i+c_i\,\omega_i\\
 &\left\{\begin{aligned}
 \overline D_i&=\overline D_i^0+c_i\,\mathrm{id}
 &\left\{\begin{aligned}&\overline D_i^0\text{ à coeff. cts}\\
                         &c_i\text{ section de }\underline{\mathcal O}_X\end{aligned}\right.\\
 (9)\quad\alpha_{\overline D_i^0}\,d\varphi&=\deleted{\unsure{\alpha_i}}\overline D_i^0\varphi
 \end{aligned}\right.
\end{aligned}}
\]}

% rewritten (8) at the foot
Considérons alors
\[
 \sum \deleted{\unsure{\alpha}}(\overline D_i)\omega_i=\overline\omega
 \tag*{(\deleted{\unsure{9}}8)}
\]
qui est encore une forme fermée. On vérifie alors \inserted{aussitôt} que l'on a
% archive sheet number: 223
